\manualchapter{Patches and troubleshooting}{sec:operation}
\subsection{Pulse-width movement}
Set Pulse 1 Duty to 0 and patch a slow 0--7\,V ramp into its duty input.
The width walks through eight steps. Leave Pulse 2's input empty to share
the ramp, then offset its knob for a different sequence of widths. Lower
both levels if the normalled sum clips.

For a stable saw bass, keep its Level fixed near the initialized value and
use an external VCA. Modulating the internal Level intentionally changes
the staircase and can produce abrupt timbral changes. If the saw repeatedly
restarts, disconnect SYNC while diagnosing the pitch.

If a fine-tuning setting differs from another chip module, remember that
Ramp VRC6's FM scale is $V/10$. A Level CV of zero silences the corresponding
voice; negative CV also clips to silence rather than inverting it.

\subsection{Cable channels, outputs and saved patches}
The widest input selects 1--16 polyphonic chip instances. Voice columns are
oscillators within each instance, not separate cable channels. Inputs read
the matching channel; supply matching pitch and level channel counts rather
than expecting a mono envelope to repeat across a polyphonic pitch cable.

Each output adds the preceding output when that earlier socket is unpatched.
The accumulated signal clips at approximately $\pm5$\,V. To hear the complete
mix, use the last output and leave earlier outputs empty. Patching an earlier
output removes that signal from the following mix. Lower source levels if
the sum clips.

Rack saves the panel parameters. Reset initializes the chip; oscillator
phase is not preserved by a patch save. Pitch/FM update each sample, while
level and other register controls update every 16 samples.
